\chapter{Dynamique relationnelle des référentiels : auto-inertie, ambivalence et principe
de Mach}
\label{chap:autoinercia-mach}

L'inertie classique caractérise la persistance d'un état de mouvement.
La question de Mach \cite{MachMechanics} déplace l'accent du corps isolé vers la
relation entre ce corps et la totalité matérielle. Dans le langage
généalogique, la même question se formule sans séparer l'observateur de la
frontière qu'il lit : un référentiel local est la restriction d'un état compatible
global, et la réponse inertielle est machienne lorsqu'elle dépend de la totalité à
travers une trace de frontière bien définie.

Le mot \emph{auto-inertiel} ne désigne pas un système déconnecté du reste.
Il désigne une section qui transporte de manière cohérente le référentiel induit par sa
propre insertion dans la totalité. Sa persistance est interne à l'état global
et se vérifie au moyen d'une connexion.

\section{Totalité, trace et réponse locale}

Soit \(\Omega_{\rm tot}\) l'espace des sections globales et
\(b:\Omega_{\rm tot}\to\mathcal B\) sa trace sur une frontière. Pour une
région ou un sommet \(v\), soit
\[
 \mathcal I_v:\Omega_{\rm tot}\longrightarrow\mathcal Y_v
\]
la réponse inertielle locale.

\begin{definition}[Factorisation machienne]
La réponse \(\mathcal I_v\) est machienne par rapport à \(b\) lorsqu'il existe
une application induite
\(\overline{\mathcal I}_v:\operatorname{im}(b)\to\mathcal Y_v\) telle que
\begin{equation}
 \mathcal I_v=\overline{\mathcal I}_v\circ b.
 \label{eq:factorizacion-machiana}
\end{equation}
\end{definition}

\begin{theorem}[Critère de factorisation par totalité]
\label{thm:criterio-mach}
Il existe une unique application
\(\overline{\mathcal I}_v:\operatorname{im}(b)\to\mathcal Y_v\) qui satisfait
\eqref{eq:factorizacion-machiana} si et seulement si \(\mathcal I_v\) est constante
sur chaque fibre de \(b\) :
\[
 b(x)=b(y)\quad\Longrightarrow\quad
 \mathcal I_v(x)=\mathcal I_v(y).
\]
Si \(b\) est surjective, \(\operatorname{im}(b)=\mathcal B\) et l'application
factorisée est définie de manière unique sur tout \(\mathcal B\).
\end{theorem}

\begin{proof}
La nécessité est immédiate. Pour la suffisance, étant donné
\(\beta\in\operatorname{im}(b)\), nous définissons
\(\overline{\mathcal I}_v(\beta)=\mathcal I_v(x)\) pour n'importe quel
\(x\in b^{-1}(\beta)\). La constance sur les fibres garantit que la
définition ne dépend pas du représentant ; l'égalité de factorisation elle-même
impose sa valeur en chaque \(\beta\in\operatorname{im}(b)\) et garantit donc
l'unicité. Si \(b\) est surjective, ce domaine coïncide avec
\(\mathcal B\).
\end{proof}

Le critère donne une formulation positive du principe de Mach : la frontière
globale contient exactement l'information de la totalité qui gouverne la
réponse locale. L'état intérieur peut conserver davantage de généalogie ;
l'inertie dépend de la classe que la trace publie.

\section{Sections auto-inertielles comme géodésiques de l'état conjoint du système, de l'observateur et de la mémoire : \(D_\gamma s_\gamma=0\)}

Soit \(\mathcal E\to\Gamma\) un fibré vectoriel d'états locaux le long
d'un chemin \(\gamma\), et soit \(D\) une connexion linéaire induite par le
transport TPK.

\begin{definition}[Section auto-inertielle]
Une section persistante \(s_\gamma\) est auto-inertielle sur un tronçon si
\begin{equation}
 D_\gamma s_\gamma=0.
 \label{eq:autoinercia-paralela}
\end{equation}
De manière équivalente et discrète, l'état à chaque sommet est le transport
parallèle de l'état précédent par le morphisme du chemin.
\end{definition}

L'équation \eqref{eq:autoinercia-paralela} conserve orientation, mémoire et
signature au sens déterminé par la connexion. Une accélération en coordonnées
peut coexister avec le transport parallèle lorsque le référentiel change
aussi. La distinction entre mouvement inertiel et force est ainsi liée à
la connexion, non à l'apparence de la trajectoire dans une carte isolée.

\begin{proposition}[Persistance et holonomie]
\label{prop:autoinercia-holonomia}
Soit \(\gamma\) un lacet et \(\Hol_D(\gamma)\) son holonomie. Une section
auto-inertielle revient exactement à son état initial si et seulement si elle appartient au
sous-espace fixe de \(\Hol_D(\gamma)\). Si la projection visible revient et que
l'état initial se trouve hors de
\(\operatorname{Fix}(\Hol_D(\gamma))\), la coordonnée revient et
l'état change.
\end{proposition}

Ce résultat relie auto-inertie et monodromie. La périodicité observable
n'implique pas la clôture du relèvement ; le déficit ou l'excès d'holonomie enregistre la
réponse accumulée du référentiel.

\section{Principe d'Ambivalence : involution bidirectionnelle des lecteurs π/φ² et φ/π²}

La cellule HMT admet deux lectures complémentaires : aréale et volumétrique. Soient
\(q_A\) le lecteur de frontière ou d'aire et \(q_V\) le lecteur d'intérieur ou de
volume. Le transport \(T_{AV}\) relie les deux sans les identifier.

\begin{principle}[Ambivalence aréale--volumétrique]
Une réalisation physique complète conserve la paire
\begin{equation}
 \bigl(q_A(x),q_V(x),T_{AV}(x)\bigr).
 \label{eq:ambivalencia-av}
\end{equation}
La lecture aréale exprime comment l'état intervient sur une frontière ; la lecture
volumétrique exprime comment il occupe et transporte l'intérieur. La troisième donnée
conserve la transformation entre les deux.
\end{principle}

Dans la mécanique dimensionnelle, la relation aréale--volumétrique organise
l'extension HMT du principe d'équivalence. Une même section peut être lue
comme réponse de frontière et comme densité intérieure. Lorsque les deux lectures
sont transportées par la même connexion locale, les descriptions inertielle et
gravitationnelle appartiennent à une seule classe d'état.

\begin{criterion}[Équivalence locale typée]
L'identification d'une réponse inertielle à une réponse gravitationnelle
est démontrée dans une région \(U\) lorsque :
\begin{enumerate}
 \item \(q_A\) et \(q_V\) sont définis sur la même section ;
 \item il existe un entrelaceur \(T_{AV}\) qui conserve leurs unités et
       orientations ;
 \item la connexion qui transporte les deux lecteurs coïncide dans \(U\) ;
 \item la différence observationnelle s'annule dans le référentiel local déclaré.
\end{enumerate}
\end{criterion}

La formulation préserve l'équivalence et montre quelles données la réalisent. Son
application à l'électron appartient à la première clôture physique de la loi des
masses ; la présente annexe utilise uniquement le schéma géométrique.

\section{Coordination causale entre transport inertiel, double lecture aréale–volumétrique et factorisation machienne}

Les trois principes s'ordonnent causalement :
\begin{equation}
 \boxed{
 \begin{aligned}
 &\text{persistance de la section}
   &&\longrightarrow&& D_\gamma s=0,\\
 &\text{double lecture du même état}
   &&\longrightarrow&& (q_A,q_V,T_{AV}),\\
 &\text{dépendance de la réponse locale à l'égard de la totalité}
   &&\longrightarrow&& \mathcal I_v=\overline{\mathcal I}_v\circ b.
 \end{aligned}}
 \label{eq:triptico-inercia-mach}
\end{equation}
Le principe d'inertie fixe le transport ; l'Ambivalence fixe les deux
lectures physiques ; Mach fixe le quotient global qui gouverne le référentiel local.
Ensemble, ils produisent une théorie relationnelle du mouvement dans laquelle l'observateur,
l'appareil et le corps appartiennent à une section commune.

\section{Mesure du mouvement et couplage des référentiels}

Mesurer une vitesse ou une accélération exige de comparer le système à une horloge et à une
règle. Ces objets font partie de l'appareil. Soient \(F_S\) le référentiel du système
et \(F_M\) le référentiel de l'observateur. Le couplage de mesure construit une
transformation relative
\[
 G_{MS}=F_M^{-1}F_S.
\]
La vitesse, l'accélération et la phase publiées sont des invariants ou des
composantes de \(G_{MS}\), non des propriétés sans référence. La répétition de
l'expérience conserve le contexte en transportant conjointement l'appareil et la
frontière.

Cette formulation retrouve l'intuition de Weyl \cite{Weyl1927} : mesurer une amplitude ou une
longueur implique de coupler une règle de jauge. HMT ajoute la généalogie de cette
règle. Le résultat peut être retracé jusqu'à l'état de l'appareil, au lecteur et au
transport qui ont fixé son échelle.

\section{Onde bidirectionnelle et orientation temporelle}

Un chemin persistant admet les canaux avancé et retardé. L'inversion
d'orientation agit sur le mot et le ledger ; dans la réalisation linéaire,
elle échange les propagateurs correspondants lorsque le Hamiltonien possède la
symétrie déclarée. Le mouvement bidirectionnel s'exprime au moyen d'une paire
\[
 (U_{\rm ret},U_{\rm adv})
\]
reliée par conjugaison et réversion temporelle.

Le ruban particule--antiparticule réalise cette dualité comme deux orientations
d'une même section. La relation dynamique peut être exprimée de façon exacte.

\begin{proposition}[Conjugaison d'orientation et réversion du propagateur]
Soit \(C\) une involution orthogonale sur la réalification quantique telle que
\[
 C J_E C^{-1}=-J_E,
 \qquad
 C H C^{-1}=H.
\]
Alors, pour
\(U(t)=\exp(-J_EHt/\hbar_{\rm int})\),
\[
 C U(t)C^{-1}=U(-t).
\]
\end{proposition}

\begin{proof}
La conjugaison commute avec la série exponentielle et transforme le générateur
\(-J_EH\) en \(+J_EH\). Le résultat est l'exponentielle évaluée en
\(-t\).
\end{proof}

L'involution de mot \(C_M=-\operatorname{rev}\) fournit la donnée
généalogique d'orientation. Lorsque sa réalisation satisfait les deux
identités de la proposition, les branches particule et antiparticule sont
représentées au moyen de propagateurs de sens temporels conjugués. La
géométrie de ruban fixe l'orientation et le Hamiltonien fixe la dynamique.

\section{Expansion globale et équation d'autocohérence}

Une cosmologie généalogique peut être formulée au moyen d'un opérateur
\(\mathcal T\) qui envoie une totalité de frontière sur la suivante :
\[
 B_{n+1}=\mathcal T(B_n).
\]
La réponse locale induite est
\[
 \mathcal I_{v,n}=\overline{\mathcal I}_v(B_n).
\]
Un régime auto-inertiel global correspond à une orbite ou à un point fixe stable
de l'opérateur de transport, tandis que l'expansion est enregistrée au moyen d'un
cocycle d'échelle \(a_{n+1}/a_n\). Cette formulation sépare trois questions :
l'existence de la dynamique, la stabilité de la solution et l'identification du
cocycle au facteur cosmologique observé.

\begin{resultbox}
L'inertie locale peut être typée comme transport parallèle d'une section ; sa
dépendance machienne s'exprime comme factorisation par la trace globale de
frontière. Le principe d'Ambivalence relie les lectures aréale et volumétrique
du même état. L'observateur intervient par le couplage des référentiels qui
rend la mesure possible.
\end{resultbox}

\begin{statusbox}
La définition de l'auto-inertie et la factorisation machienne sont des résultats
récupérés. Le critère de dimension et l'instrument de mesure proviennent des
chapitres précédents. La synthèse \eqref{eq:triptico-inercia-mach} est une
nouvelle formalisation d'une architecture autorale préexistante. Une équation
cosmologique quantitative exige de fixer \(\mathcal T\), sa métrique et ses données
de frontière.
\end{statusbox}

L'attribution machienne est caractérisée par la constance de la réponse
locale sur les fibres de la trace globale. L'auto-inertie est caractérisée
par le transport parallèle et par le sous-espace fixe de l'holonomie. Dans une
réalisation cosmologique, l'opérateur \(\mathcal T\), sa métrique et le cocycle
d'échelle transforment ces conditions structurelles en une évolution
quantitative de la totalité.
